\section{Ý Tưởng Đề Tài}

\subsection{Bối Cảnh Thực Tiễn và Bài toán}
\contrib

Trong những năm trở lại đây, du lịch Việt Nam đã trở thành một đầu tàu mũi nhọn
đóng vai trò rất lớn đối với nền kinh tế Việt Nam. Nó không chỉ  sở hữu tài nguyên phong phú và đa dạng,
mà còn có rất nhiều các di tích lịch sử ngàn năm văn hiến, danh lam thắng cảnh kỳ
vĩ đến các di sản văn hóa phi vật thể và ẩm thực đặc sắc của
từng vùng miền. Trong những năm gần đây, đặc biệt là giai đoạn phục
hồi mạnh mẽ sau đại dịch, ngành du lịch đã ghi nhận những bước chuyển
mình vượt bậc với các số liệu tăng trưởng ấn tượng. Theo số
liệu chính thức từ Tổng cục Thống kê, vào năm 2025, Việt Nam đã đón
gần 21.2 triệu lượt khách quốc tế, đạt mức tăng trưởng hơn 20\% so
với năm 2024 và thiết lập kỷ lục lịch sử mới về lượng khách ngoại
quốc ghé thăm trong một năm~\cite{gso2025}. Bên cạnh đó, thị trường
du lịch nội địa cũng ghi nhận sức phát triển sôi động khi
phục vụ khoảng 135.5 triệu lượt khách, từ đó nâng tổng thu từ khách
du lịch lần đầu tiên trong lịch sử vượt mốc 1 triệu tỷ đồng~\cite{vnat2025}.
Nhằm tiếp tục phát huy đà tăng trưởng này, ngành du lịch nước nhà đã
đặt mục tiêu lớn cho năm 2026 với việc đón 25 triệu lượt khách quốc
tế và đạt tổng doanh thu ước tính khoảng 1.125 triệu tỷ đồng~\cite{vnat2026},
khẳng định vai trò cốt lõi của du lịch như một ngành kinh tế mũi
nhọn của đất nước. Tuy nhiên, dù cho với những số liệu tích cực đó, ngành du lịch
nước vẫn còn nhiều hạn chế dẫn đến nhiều bất cập cho du khách trong nước và ngoài nước
khi đến trải nghiệm.

\textbf{Trước hết} là nước ta đang đối với mặt với một thực trạng là tỷ lệ khách quay lại khá thấp. Khi chúng ta
so sánh với các quốc gia trong khu vực như Thái Lan --- nơi có tỷ lệ
khách quay lại đạt đến gần 80\% thì tại Việt Nam, tỷ lệ này chỉ dao
động từ 10\% đến 40\%~\cite{pata2020}. Điều này phản ánh thực tế rằng
các sản phẩm du lịch hiện nay đang thiếu đi một cái gì đó để níu kéo khác hàng, khiến du khách thường chỉ tới một
lần duy nhất chỉ để ngắm cảnh mà không thấy được lý do mạnh mẽ để quay
lại thêm một lần nữa. Đặc biệt quan trọng là công tác quản lý điểm
đến tại nhiều địa phương bộc lộ ra nhiều hạn chế, chưa loại bỏ được
tình trạng đeo bám, ép giá, hay vấn đề vệ sinh môi trường, qua đó
để lại những ấn tượng xấu trong lòng bạn bè quốc tế.

Tiếp đó chính là việc thông tin và dữ liệu phân tán khắp
mọi nơi trong thời đại số. Hiện nay, dữ liệu về các điểm đến, hoạt
động văn hóa nghệ thuật, hay ẩm thực đặc sản phân tán ra trên nhiều
nền tảng số như Google Maps, các trang mạng xã hội, các blog cá nhân
hay các ứng dụng cục bộ của từng địa phương. Điều này buộc du khách
phải tiến hành công việc tìm kiếm và đối chiếu thông tin khá thủ
công trước và trong suốt chuyến đi. Mặt cách, công cụ số hóa hiện nay
hầu như chỉ cung cấp thông tin mặt ngoài như tọa độ, thời gian đóng
mở cửa và đánh giá sao chung chung, hoàn toàn thiếu đi khả năng giải
nghĩa các câu chuyện lịch sử, giá trị nhân văn và ý niệm văn hóa sâu
xa đằng sau các di tích cổ kính.

Và cuối cùng là việc thiếu hụt các giải pháp công nghệ tự
động hóa tùy theo điều kiện thực tế. Trải nghiệm của du khách tại
điểm đến bị chi phối khá nhiều bởi các yếu tố ngoại cảnh thay đổi
liên tục như thời tiết, tình trạng ùn tắc giao thông, khoảng cách địa
lý hay mật độ đông đúc. Một di sản lịch sử có giá trị văn hóa lớn,
nhưng nếu du khách tới tham quan vào thời điểm mưa bão hay quá tải
chen chúc thì trải nghiệm họ thu nhận được sẽ kém hiệu quả đi nhiều. Rõ ràng, việc các hệ thống hiện nay không thể
tự động điều chỉnh đề xuất tùy theo bối cảnh môi trường thời gian
thực đã tạo ra khoảng cách rất lớn giữa kỳ vọng của du khách và
khả năng đáp ứng của dịch vụ.

Từ những khó khăn thực tế nêu trên, ta đặt ra câu hỏi rằng làm sao
ta có thể phát triển một giải pháp công nghệ có khả năng tổng hợp các
nguồn dữ liệu phân tán thành một hệ thống chung, đồng thời tự
động hóa việc đưa ra quyết định hỗ trợ du khách tùy theo điều kiện
ngoại cảnh thay đổi liên tục. Điều này tạo ra nhu cầu cần thiết cho
một mô hình hệ thống có khả năng thích ứng bối cảnh, xử lý dữ liệu
lớn về văn hóa bản địa và dẫn dắt hành trình du lịch của người dùng
một cách hiệu quả.

\subsection{Problem Statement}
\contrib
Trước hết nếu ta nhìn theo góc độ của người dùng, đây là một hệ thống thông minh hỗ trợ khách du lịch
tìm kiếm địa điểm du lịch thú vị, nhưng song với đó cũng cung cấp thêm nhiều kiến thức về văn hóa của nơi đó
nhằm tạo nên một sự hứng thú cho họ. Còn nếu nhìn từ góc độ
kỹ thuật, thì lại là một bài toán tối ưu hóa quyết định với nhiều mục tiêu
và ràng buộc thay đổi theo thời gian thực, đòi hỏi ta phải mô hình hóa được
sở thích của người dùng từ thực thể đời thường, thành một đối tượng toán học và từ đó
có thể cá nhân hóa theo sở thích, tương tác của họ.
Đây chính là động lực chính để phát triển hệ thống
\textbf{BeVietnam}, một giải pháp du lịch thông minh dựa trên mô hình
\textbf{agent-based system} nhằm giảm tình trạng phân tán thông
tin, hỗ trợ đề xuất linh hoạt tùy theo điều kiện thực tế và dẫn dắt hành
trình trải nghiệm thông qua các câu chuyện. Và để làm được điều đó, nhóm đã đề ra
ba tính năng chính và là mục tiêu cốt lõi của dự án như sau:

\textit{Đầu tiên}, phần lớn các du khách khi đến một điểm du lịch nào đó thì học chủ yếu để ngắm cảnh 
và vui chơi. Tuy nhiên, nếu một địa điểm đó chỉ được dùng để ngắm xong rồi về thì nó sẽ tạo cho họ
một cảm giác chán và thiếu hụt đi cái gì đó. Họ chỉ ngắm nhưng họ không hiểu thật sự phía sâu bên trong mỗi
cảnh đẹp, mỗi kỳ quan kia đang chất chứa điều gì ở đó. Chính vì lẽ đó để giảm bớt vấn đề cảnh quan, địa điểm thiếu đi chiều sâu dẫn
đến tỷ lệ quay lại thấp, nhóm đã nghiên cứu và xây dựng tính năng nhiệm vụ văn hóa
storyline theo phong cách geocaching thực địa. Tại nơi đây, người dùng sẽ đóng vai trò như là một \textbf{nhà lữ hành}
thực thụ, họ sẽ được giao cho những nhiệm vụ xoay quanh các kiến trúc cổ, các danh lam thắng cảnh, các địa vui chơi mới và sau đó
họ sẽ tiến hành chụp ảnh, checkin nộp lên hệ thống để kiểm chứng và hoàn thành nhiệm vụ. Hệ thống tổ chức toàn bộ
nhiệm vụ thành một \textbf{Mission Path} trực quan để giữ chân người dùng
trong đó mỗi nút trên lộ trình là một nhiệm vụ văn hóa được
nạp trực tiếp từ kho câu hỏi đã được xây dựng từ trước từ rất nhiều cuốn sách văn hóa về điểm đó. Thay
vì chỉ ngắm cảnh thụ động, du khách buộc phải di chuyển đến tọa độ GPS
của điểm đến và chụp ảnh minh chứng các yếu tố lịch sử, kiến trúc. Ảnh
minh chứng này được một mô hình thị giác chấm
điểm theo một thang đánh giá để quyết định mở khóa nút tiếp
theo, qua đó biến chuyến đi thành hành trình tương tác rõ ràng hơn, ly kỳ hơn và lý thú hơn.

\textit{Tiếp theo}, nhằm triệt tiêu vấn đề phân tán thông tin du lịch, nhóm phát
triển tính năng bảng tin gợi ý feed thông minh. Tính năng này hoạt động
giống như một bộ lọc trung tâm, tự động tổng hợp nguồn tri thức văn
hóa chính thống từ nhiều cuốn sách và trang web khác nhau để trả về các bài viết giới thiệu điểm đến trực quan 
mới mẻ và thú vị, giúp du khách nắm bắt trọn vẹn giá trị di sản mà không phải tìm
kiếm thủ công trên nhiều nền tảng khác nhau. Khi phát triển tính năng này, nhóm làm ra với mong muốn là trước khi họ đến
một nơi nào đó, họ sẽ có thể tìm hiểu trước về địa điểm đó, biết được địa điểm sắp tới của mình là gì, biết được tại sao nó lại nhiều
người đến thăm, biết tại sao nó nổi tiếng, từ đó mà có thể tọa nên một sự tò mò và kích thích từ sâu bên trong họ.

\textit{Cuối cùng}, để khắc phục điểm yếu về việc phản ứng theo điều kiện ngoại
cảnh thực tế, nhóm thiết kế tính năng bubble map tương tác hiển thị các
POI thời gian thực quanh vị trí người dùng. Hệ thống lấy vị trí GPS
thực của du khách, truy vấn các địa điểm trong bán kính 3km qua nguồn
dữ liệu \textbf{Foursquare}, đồng thời thu thập thời tiết khu vực hiện tại (nhiệt
độ, chỉ số UV ước lượng, lượng mưa) bằng \textbf{OpenWeather API}. 
Kích thước mỗi bong bóng co giãn theo khoảng cách thực tế từ địa điểm tới người dùng, 
qua đó giúp du khách nhận biết nhanh nơi nào đáng cân nhắc nhất tại thời điểm hiện tại.
Điều này làm ra bởi lẽ, khi du khách muốn đi một nơi nào đó họ sẽ rất muốn biết rằng nơi đó thời tiết như nào có đang mưa
hay không, có đông hay không, có nhiều review đánh giá tốt hay không? Và đặc biệt là làm sao để họ có thể tiếp cận được những điều đó.
Ví dụ như họ vô Google Maps thì làm thế nào mà họ có thể biết được quán này muốn đến, quán này nên đi, điểm kia nên tới một cách tiết kiệm nhất có thể,
họ không thể nào mà dành hàng giờ đồng hồ để kiểm tra, tìm hiểu về nó xong rồi khi sắp tới nơi thì lại mắc mưa.

\subsection{Tầm Nhìn Sản Phẩm}
\contrib

\textit{BeVietnam} được định hướng phát triển thành một hệ sinh thái du lịch văn
hóa thông minh, đầu tiên sẽ tập trung thí điểm tại Huế. Thực ra ở giai
đoạn đầu nhóm từng muốn bao phủ nhiều tỉnh thành cùng lúc, nhưng khi bắt
tay chuẩn bị dữ liệu văn hóa thì thấy nội dung quá dàn trải và khó kiểm
chứng nguồn, nên nhóm chủ động thu hẹp lại còn điểm duy nhất là Huế để giữ
được chất lượng dữ liệu và nội dung. Bởi lẽ khi thực hiện điều này, nhóm phải
bỏ thời gian công sức rất nhiều để tìm kiếm các nguồn tài liệu chính thống, các cuốn sách hay
về điểm đó, sau đó sẽ tiến hành quá trình ingest và tạo câu hỏi, đòi hỏi một khoảng thời gian dài và
nhiều công sức. Chính vì lẽ đó trong khuôn khổ đồ án này nhóm chỉ tập trung vào tỉnh thành đó chính là Huế.

Sản phẩm hướng đến các câu hỏi cho khách du lịch là \textit{nên đi đâu vào lúc này, vì sao địa điểm đó
đáng đi và khi đến đó thì nên làm gì} thông qua đó mà hệ thống có thể giúp họ có câu trả lời cho riêng mình.
Từ tầm nhìn cốt lõi đó, sản phẩm được triển khai vói ba trải nghiệm
tương tác chính như sau: \textbf{Một là, bảng tin gợi ý địa điểm tương tự một feed
mạng xã hội năng động}, trong đó mỗi đề xuất được tính toán và sắp xếp
theo sự kết hợp giữa vị trí của người dùng, thời tiết hiện tại, điều
kiện giao thông thực tế và sở thích cá nhân. \textbf{Hai là, bubble map hiển
thị trực quan các điểm} đến với kích thước bong bóng thay đổi theo mức
độ phù hợp của từng địa điểm ở thời điểm hiện tại. \textbf{Ba là, hệ thống
nhiệm vụ văn hóa} theo phong cách geocaching, buộc người dùng phải trực
tiếp di chuyển đến thực địa, check-in và chụp ảnh minh chứng để mở khóa
các nội dung giải nghĩa văn hóa sâu sắc tiếp theo.

\subsection{Đối Tượng Khách Hàng Mục Tiêu}
\contrib

Đối tượng người dùng mục tiêu của hệ thống BeVietnam bao gồm nhiều loại khách hàng khác nhau từ trong
nước tới ngoài nước. Trước hết, nhóm du khách nội địa và quốc tế có nhu
cầu tìm hiểu sâu sắc về lịch sử, kiến trúc cổ kính và ẩm thực độc đáo
tại địa phương nhưng thường gặp rào cản về mặt ngôn ngữ hoặc thiếu thốn
thông tin giải nghĩa văn hóa có hệ thống. Kế đến, nhóm người đam mê
du lịch tự túc và mong muốn tự chủ hoàn toàn lịch trình của mình, rất
cần các đề xuất cá nhân hóa linh hoạt theo thời gian thực thay vì tuân
theo các tour tuyến cố định và nhàm chán. Sau cùng, nhóm người dùng trẻ
yêu thích khám phá công nghệ mới và các trải nghiệm trò chơi hóa, mong
muốn tiếp thu kiến thức văn hóa một cách tự nhiên, chủ động thông qua
việc hoàn thành các thử thách thực địa đầy thú vị.
